%--------------------
% Packages
% -------------------
\documentclass[11pt,a4paper]{article}
\usepackage[utf8x]{inputenc}
\usepackage[T1]{fontenc}
%\usepackage{gentium}
\usepackage{mathptmx} % Use Times Font


\usepackage[pdftex]{graphicx} % Required for including pictures
\usepackage[english]{babel} % English translations
\usepackage[pdftex,linkcolor=black,pdfborder={0 0 0}]{hyperref} % Format links for pdf
\usepackage{calc} % To reset the counter in the document after title page
\usepackage{enumitem} % Includes lists
\usepackage{ctable}
\usepackage{subcaption}

\frenchspacing % No double spacing between sentences
\linespread{1.0} % Set linespace
\usepackage[a4paper, lmargin=0.1111\paperwidth, rmargin=0.1111\paperwidth, tmargin=0.1111\paperheight, bmargin=0.1111\paperheight]{geometry} %margins
%\usepackage{parskip}

\usepackage[all]{nowidow} % Tries to remove widows
\usepackage[protrusion=true,expansion=true]{microtype} % Improves typography, load after fontpackage is selected

\usepackage{lipsum} % Used for inserting dummy 'Lorem ipsum' text into the template


%-----------------------
% Set pdf information and add title, fill in the fields
%-----------------------
\hypersetup{ 	
pdfsubject = {PhD candidate report, mid-term evaluation},
pdftitle = {},
pdfauthor = {}
}

\title{\textbf{PhD candidate report, mid-term evaluation}\\
\textit{}}
\author{}
\date{}

%-----------------------
% Begin document
%-----------------------
\begin{document} %All text i dokumentet hamnar mellan dessa taggar, allt ovanför är formatering av dokumentet

\begin{figure}
    \centering
    \includegraphics[width=0.5\textwidth]{uialogo.jpg}
\end{figure}
\maketitle

\begin{flushleft}

\begin{table}[h]
    \begin{tabular}{l l}
    &\\
    &\\
    &\\
    &\\
    &\\
    &\\
    &\\
    &\\
    &\\
         PhD student:& \textit{}	\\
         &\\
        University of Agder (UiA) & \\
        Faculty of Engineering and Science& \\
        Department of ICT & \\
        & \\
        Postal address:& \\
        Post Box 422& \\
        NO-4604 Kristiansand& \\
        Norway& \\
        & \\
        Visiting address: &	\textit{Jon Lilletuns vei 9, Grimstad, Norway} \\
        
        Office phone:&	\textit{} \\
        Mobile phone:&	\textit{} \\
        Email:&	\textit{} \\
    \end{tabular}
\end{table}
    
\end{flushleft}

\textbf{Information}
    \begin{table}[h]
        \centering
        \begin{tabular}{l p{9cm}}
PhD programme:&	\textit{PhD programme in Engineering and Science, Specialisation in ICT, Artificial Intelligence} \\
Project title:&	\textit{} \\
PhD Candidate:&	\textit{} \\
Principal supervisor:&	\textit{} \\
Co-supervisor:&	\textit{} \\
Co-supervisor:&	\textit{} \\
Start date of the PhD project:&	\textit{} \\
Expected date of completion:&	\textit{} \\

        \end{tabular}
    \end{table}

\textbf{PhD Candidate}

\begin{tabular}{p{0.2\textwidth}cp{0.35\textwidth}cp{0.25\textwidth}}
    \\
    &&&&\\
    \\
  \cline{1-1} \cline{3-3} \cline{5-5} \\
  Date && Signature, PhD Student && Printed Name \\
  \\
  \specialrule{.2em}{.1em}{.1em}
  \hline
  \\
  \\
\end{tabular}

The PhD candidate must submit the mid-term report to the evaluators 3 weeks prior to the date of 

evaluation. 

\begin{tabular}{@{}p{5cm}p{5cm}@{}}
\\
\textbf{Date} for mid-term evaluation: & \hrulefill \\
\\
\\
\end{tabular}

\large{\textbf{Oral Presentation}}
\begin{itemize}
    \item 30-45 minutes presentation by the PhD candidate of key elements from the report.
    \item Discussion of report and presentation, including discussion of the wider thematic framework.
    \item Discussion of plans for completion of project and studies.

\end{itemize}

\newpage

\section{Problem Statement}


\section{Introduction}


\subsection{Topic 1}


\subsection{Topic 2}


\subsection{Topic 3}


\section{Research Questions Addressed}


\subsection{RQ1}


\subsection{RQ2}


\section{Remaining Research Questions and Methodology}

\subsection{RQ3}


\subsection{RQ4}


\section{Progress Plan}
\begin{enumerate}    
  \item Year 1

  \begin{enumerate}
    \item Quarter 1
    
    \begin{enumerate}
      \item
    \end{enumerate}
    
    \item Quarter 2
    
    \begin{enumerate} 
      \item 
    \end{enumerate}
    \item Quarter 3
    
    \begin{enumerate}
      \item 
    \end{enumerate}
    
    \item Quarter 4
    
    \begin{enumerate}
      \item 
    \end{enumerate}


  \end{enumerate}

\item Year 2

  \begin{enumerate}
    
    \item Quarter 1
    
    \begin{enumerate}
      \item 
    \end{enumerate}
    
    \item Quarter 2
    
    \begin{enumerate}
      \item 
    \end{enumerate}
    
    \item Quarter 3
    
    \begin{enumerate}
      \item 
    \end{enumerate}
    
    \item Quarter 4
    
    \begin{enumerate}
      \item 
    \end{enumerate}

  \end{enumerate}

\item Year 3

  \begin{enumerate}

    \item Quarter 1
    
    \begin{enumerate}
      \item 
    \end{enumerate}
    
    \item Quarter 2
    
    \begin{enumerate}
      \item 
    \end{enumerate}
    \item Quarter 3
    
    \begin{enumerate}
      \item 
    \end{enumerate}
    
    \item Quarter 4
    
    \begin{enumerate}
      \item 
    \end{enumerate}

  \end{enumerate}

\end{enumerate}

\subsection{Publications}
    \begin{table}[h]
        \centering
        \begin{tabular}{|p{6cm}|p{7cm}|p{2cm}|}
        \hline
            Journal/Conference & Title & Status\\
            \hline
             &  &  \\
            \hline
             &  &  \\
            \hline
             &  &  \\
            \hline
        \end{tabular}
        \caption{List of publications.}
        \label{tab:publications}
    \end{table}

\section{Research Process}

\subsection{Reflections}
\begin{itemize}
    \item 
\end{itemize}
\subsection{Pointing Ahead}
\begin{figure}[h]
    \centering
    \includegraphics[width=0.8\linewidth]{Example gantt chart.png}
    \caption{Future plan.}
    \label{fig:gantt}
\end{figure}
\bibliographystyle{acm}
\bibliography{references}

\end{document}
